% =====================================================================
% INTRODUCTION
% =====================================================================
\section{Introduction}
\label{sec:introduction}

A language model in production is rarely reachable directly. It sits
behind application code that decides what enters its context, what it
may retrieve, what is done with what it emits, and what side effects
its output may cause: an input filter, a retriever over a corpus users
can write to, a decoder turning an uploaded image into prompt text, an
output filter, an agent layer that executes model-authored queries.
Each is a place where attacker-supplied data can acquire the authority
of an instruction --- the structural observation behind indirect prompt
injection~\cite{greshake2023}, since the model cannot tell an
instruction it was \emph{given} from one it was \emph{shown} when both
arrive as the same tokens in the same context. Greshake et al.\ put the
consequence as processing retrieved untrusted data being analogous to
executing arbitrary code~\cite[Sec.~2]{greshake2023}.

The security question worth asking about a deployed system is therefore
not primarily about the model. A model probed in isolation cannot
exhibit cross-tenant database reads, cannot leak a coupon a server-side
authorisation check would never have placed in its context, and cannot
be poisoned through a comment form; those failures belong to a
pipeline. Most red-teaming work nevertheless takes the model as the
unit of analysis --- reasonably, since models generalise across
deployments --- leaving a gap: an organisation running an LLM-backed
application has many model-level results and little guidance on
measuring \emph{its own system} against \emph{its own policy} in a way
another engineer can replay.

\subsection{Detection is the Hard Part}
\label{sec:intro:detection}

Closing that gap runs into the step automated LLM security testing can
neither avoid nor generalise: deciding whether an output constitutes a
failure. The decision is deployment-specific --- the same string is a
policy breach in one system and an unremarkable answer in another ---
so a portable detector can only ask a generic question (\emph{does this
look like a jailbreak?}) and inherits a precision ceiling set by how
often a benign output resembles a violating one. We measure that
ceiling rather than assume it (Sec.~\ref{sec:res:agreement}).

An assessment can do better when the application's policy is legible
enough to transcribe, the condition this study exploits. The target
names, in its own source, the exact word each persona must not emit; it
synthesises PII of exactly specified shapes; it gives each seeded
account an unguessable routing flag in an isolated database. Each
constant is ground truth in the strict sense --- its presence in an
output is a violation and its absence is not --- which converts a
subjective scoring problem into an objective one at the cost of
detectors that transfer to no other system: precisely the trade an
assessment, rather than a benchmark, should make. Two problems follow.
Ground truth must \emph{reach} the detector, since a suite sweeping
five personas in one process must score each attempt against \emph{that
attempt's} secret. And the detector must still \emph{recognise} the
violation --- a secret rewritten as \code{C-H-E-E-S-E} is still a leak
--- where every refinement that buys sensitivity spends precision.

\subsection{Contributions}
\label{sec:intro:contributions}

This paper reports a policy-aware, application-level assessment of
PwnzzAI Shop, an intentionally vulnerable LLM-backed web application,
built entirely as plugins for Garak, the open-source LLM security
scanner.\footnote{\url{https://github.com/NVIDIA/garak}. Named as the
tool the assessment is expressed in, not as a source of claims.} The
suite, its artefacts and its documentation are the companion
record~\cite{pwnzzrepo,pwnzzdocsoverview}.

\begin{enumerate}
\item \textbf{A Garak-native architecture for assessing an application
rather than a model} (Sec.~\ref{sec:architecture}): nine surfaces
wrapped as generators and registered into Garak's namespace and plugin
cache at run time without copying files into the installation, so Garak
core is unmodified and every measurement replays through the stock CLI.

\item \textbf{Ground-truth detection on a response side-channel}
(Sec.~\ref{sec:arch:notes}): twelve detectors testing a specific
violation against constants transcribed from the pinned source and
re-verified by contract tests, with per-attempt ground truth travelling
in \code{Message.\bk notes}, and abstention --- \code{None}, never
$0.0$ --- whenever a detector cannot judge. A calibrated
LLM-as-a-judge (Sec.~\ref{sec:arch:judge}) covers what ground truth
cannot see, a response satisfying the attacker without matching the
pattern; it is off by default, never primary, and its schema-order
calibration is reported rather than assumed.

\item \textbf{A layer-isolating design with built-in controls}
(Sec.~\ref{sec:methodology}): one attack corpus swept across five
persona levels and ten guardrail configurations differing only in which
defensive layer is present; poisoning measured against a paired clean
control fitted in the same run; a retrieval mitigation toggled with
corpus and prompts fixed; and a delivery-integrity check excluding
indirect attempts whose payload never arrived.

\item \textbf{Results over 678 evaluated attempts}
(Sec.~\ref{sec:results}): success concentrated in techniques that never
name their objective; two of ten guardrails no better than none; a
poisoning threshold lagging the effect by four budget steps; a generic
detector at perfect recall and $0.293$ precision.

\item \textbf{An ablation study} (Sec.~\ref{sec:ablation}) separating
each defensive layer from each element of the measurement apparatus: of
ten components, three are guarantees these runs never exercised, five
change a reported number --- one by $3.4\times$ --- and two abolish a
measurement rather than degrading it.
Sec.~\ref{sec:disc:mitigations} then ties each mitigation to its
supporting detector signal and the layer it acts on, with two
false-positive classes reported against our own detectors.
\end{enumerate}
